\CNsection{}{Introduction}

\textbf{Advanced Communication Networks} is a 3-credit course (48 hours, 16 weeks \ensuremath{\times} 3 hours) for senior undergraduate and graduate students of electrical and telecommunication engineering, taught as lectures with laboratories and a project. The course is a coherent technical journey through the evolution of mobile networks from GSM (2G) to LTE (4G), 5G NR/\allowbreak{}SA and the emerging concepts of 6G — not a collection of separate technologies, but a continuous analysis of \textbf{why} each generation appeared, \textbf{what} problem it solved and \textbf{how} traffic flows through its architecture. The course is built on a “why /\allowbreak{} what /\allowbreak{} how” logic so that students can analyse and evaluate architectural decisions instead of memorising lists of protocols.

\CNsection{}{Structure and Content}

The material is organised in \textbf{26 modules}:

\begin{itemize}
\tightlist
\item
  \textbf{2G/\allowbreak{}3G foundations (modules 1–4):} GSM architecture, GPRS and EDGE, circuit switching and SS7 signalling, 2G/\allowbreak{}3G core networks and the basics of mobility management and handover.
\item
  \textbf{4G access and core (modules 5–13):} LTE architecture, the E-UTRAN access network, the protocol stack, resource scheduling, the random access channel (RACH), the Evolved Packet Core (EPC), the Diameter protocol, the HSS and authentication, the MME, and the S-GW/\allowbreak{}P-GW with GTP tunnelling.
\item
  \textbf{Software-driven architectures (modules 15–18):} software-defined networking (SDN), network function virtualisation (NFV), the 3GPP standardisation process and the 5G service-based architecture (SBA).
\item
  \textbf{5G core and procedures (modules 14, 19–22):} the 5G core architecture, signalling procedures (registration, 5G-AKA authentication, PDU session establishment and release, paging, handover), network slicing and the evolution of the QoS framework (from QCI to 5QI).
\item
  \textbf{Planes and outlook (modules 23–26):} the user plane and the control plane, a comprehensive synthesis of network evolution, and 6G research directions including ISAC, RIS, AI-native networks, cell-free MIMO and sub-THz.
\end{itemize}

Every module comes with comparison tables, architecture diagrams, step-by-step signalling sequence diagrams and review questions. A prerequisites review recalls the basics of OSI, RF, modulation, OFDM, IP networking, probability and decibel arithmetic.

\begin{CNbox}[unbreakable]{obj}{Learning Objectives}

By the end of the course, students can trace the evolution of the generations; analyse core architectures (circuit-switched, packet-switched, the all-IP EPC and the service-based 5G core); compare mobility management and handover across generations; apply network slicing, QoS, SDN and NFV in modern networks; explain 5G core procedures at the signalling level; and critically assess 6G research directions.

\end{CNbox}

\Needspace{15\baselineskip}

\CNsection{}{Assessment and Practical Work}

{\def\LTcaptype{none} % do not increment counter
\Needspace{9\baselineskip}
\begin{longtable}[c]{@{}cc@{}}
\toprule\noalign{}
\rowcolor{CNink}\CNth{Component} & \CNth{Weight} \\
\CNheadrulesep
\endhead
\bottomrule\noalign{}
\endlastfoot
Exams (midterm 20\% + final 30\%) & 50\% \\
Laboratories (6 labs \ensuremath{\times} 5\%) & 30\% \\
Project (proposal 5\% + report and demo 15\%) & 20\% \\
\end{longtable}
}

The laboratory assignments cover the LTE link budget, EPC signalling traces, 5G NR scheduling and resource allocation, SBA service discovery, network slicing configuration and user-plane traffic tracing. The project, individual or in pairs, is defined on topics such as a comparative analysis of LTE versus 5G, an enterprise network-slicing design, implementing a scheduling algorithm with the simulator, or signalling analysis with real packet captures.

\CNsection{}{Main References}

Sesia et al. (\emph{LTE – The UMTS Long Term Evolution}), Dahlman et al. (\emph{5G NR: The Next Generation Wireless Access Technology}), Holma and Toskala (\emph{LTE for UMTS}), the 3GPP specifications (TS 23.501, 23.502, 24.501, 38.300) and Rappaport (\emph{Wireless Communications}).
